\documentclass{ermacrs2023}

\title{XI Encontro Regional de Matemática Aplicada e Computacional do Rio Grande do Sul -- Instruções para Resumos}

% Usar \author{}, com o argumento em branco, para suprimir os nomes
% \author{}
\author{Nome Sobrenome$^{1}$, \underline{Nome Sobrenome}$^{2}$, Nome Sobrenome$^{3}$, Nome Sobrenome$^{4}$}

% Usar \info{}, com o argumento em branco, para suprimir as filiações e os endereços eletrônicos
% \info{}
\info{$^1$Universidade, xxxx@yyy.zzz; \\ $^2$Universidade, aaaa@bbb.ccc; \\$^3$Universidade, nnnn@ppp.qqq; \\ $^4$Universidade, nnnn@ppp.qqq.}
  
\begin{document}
  \maketitle
  
  Os \textbf{resumos} devem ter até 2 páginas, sem contar a página de referências, ser redigidos em \textbf{língua portuguesa} e estar de acordo com este modelo modelo \LaTeX.
  
  A submissão de trabalhos deverá ocorrer pela plataforma \texttt{Even3}, em endereço eletrônico divulgado na página do evento. O trabalho deve ser submetido em formato \texttt{.pdf}. Também será solicitada a submissão de uma versão do trabalho não identificada. Não serão aceitas propostas enviadas por qualquer outro meio, tampouco após os prazos.

  No momento da submissão os trabalhos devem ser classificados em uma das sessões técnicas: ST1 – Análise Aplicada; ST2 – Biomatemática; ST3 – Física-Matemática, Mecânica dos Fluidos e Sistemas Dinâmicos; ST4 – Matemática Aplicada à Engenharia; ST5 – Modelagem Matemática e Computacional; ST6 – Métodos Estocásticos e Estatísticos; ST7 – Métodos Numéricos e Computação Científica; ST8 – Matemática Discreta e Otimização; ST9 – Ensino e Extensão em Matemática Aplicada e Computacional; ST10 – Análise; ST11 – Geometria; ST12 – Álgebra; ST13 – Teoria dos Números; ST14 – Modelagem Matemática Aplicada às Ciências Naturais.

  \vspace{0.5cm}
  {\bf Submissão de trabalhos:}

  A data final para a submissão dos trabalhos é \textbf{31/03/2023}, com divulgação dos resultados ocorrendo até o dia \textbf{30/04/2023}.
  
  Demais informações sobre o evento e sobre a submissão de trabalhos podem ser encontradas no site \href{https://wp.ufpel.edu.br/ermacrs23}{\textbf{\underline{https://wp.ufpel.edu.br/ermacrs23}}}. 
  
  Todas as equações apresentadas no texto devem ser sequencialmente numeradas, conforme o exemplo a seguir
  \begin{equation}\label{eq:equacao-quadratica}
    ax^{2} + bx + c = 0.
  \end{equation}
 
  Devido a natureza do resumo, o texto não deve ser dividido em seções, com exceção apenas de uma seção de referências e, opcionalmente, uma seção de agradecimento, ambas não numeradas. Além disso, recomenda-se evitar o uso tabelas ou figuras no corpo do texto. Se necessário, as figuras devem ser inseridas no corpo do texto através do ambiente \texttt{figure} e as tabelas pelo ambiente \texttt{table}. As figuras e as tabelas apresentadas no corpo do texto devem estar próximas do local onde são inicialmente referidas. Todas as figuras e tabelas presentes no texto deverem ser numeradas, centralizadas e ter a sua fonte adequadamente referenciada. Adicionalmente, elementos textuais presentes nas figuras devem ser legíveis, com tamanho de fonte compatível à do corpo do texto, e elas devem ter uma resolução mínima de 300 dpi. Ainda, a legenda da tabela deve estar acima dela. Como exemplos, temos a Figura~\ref{fig:logo-ermacrs} e a Tabela~\ref{tab:exemplo-tabela} a seguir.
  \begin{figure}[ht]
    \centering
    \includegraphics[width=0.3 \textwidth]{logoermacrs.jpeg}
    \caption{Exemplo de figura. \\ Fonte: do autor.} 
    \label{fig:logo-ermacrs}
  \end{figure}

  \begin{table}[ht]
   \centering
   \caption{Exemplo de tabela$^{*}$.} \label{tab:exemplo-tabela}
   \begin{tabular}{ccc}
      \toprule
        Elemento X & Elemento Y & Elemento Z \\ \midrule
        0,1 & 0,3 & 0,4 \\ \bottomrule
    \end{tabular}
    \caption*{Fonte: do autor. \\ $^{*}$Alguma informação adicional sobre a tabela, se tiver.}
  \end{table}

  Na seção de referências, devem estar listados todos as obras citadas, por ordem alfabética de sobrenomes. Para citações, os autores devem utilizar comandos como o \texttt{citep} e o \texttt{citeyear}, conforme os exemplos:: \citep{strogatz2015,EadesWormald,Brusquetta}, ``ainda Viana (\citeyear{Viana}) diz que...''.

  \textbf{Esta seção tem um limite de tamanho de uma página. Além disso, esta seção não contabiliza no número total de páginas de cada artigo.}
  
  \section*{Agradecimentos}
  
  Agradecimentos, se houver, devem vir logo após a última seção do corpo do texto e antes da seção “REFERÊNCIAS”.

  \vfill\pagebreak
  
  \begin{thebibliography}{99}
    \bibitem[BRUSQUETTA {\it et al.}, 2019]{Brusquetta}Brusqueta, M., Picotti E., Mion, E., Chen, Y., Beghi, A., Minen, D. A., ``Nonlinear Model Predictive Control based Virtual Driver for high performance driving''. Em: {\bf Annals of IEEE Conference on Control Technology and Applications} (2019), Hong Kong, China.
      
    \bibitem[EADES; WORMALD,  1994]{EadesWormald} Eades, P.,  Wormald, N. C.  ``Edge crossings in drawings of bipartite graphs''. Em: {\bf Algorithmica} 11 (1994), pp. 379--403. DOI: \texttt{10.1007/BF01187020}.

    \bibitem[XI ERMARC-RS, 2023]{Ermac} ERMAC--RS. Site oficial do XI Congresso Regional de Matemática Aplicada e Computacional do Rio Grande do Sul. Online. Acessado em 07/02/2023, \href{https://wp.ufpel.edu.br/ermacrs23}{\textbf{\underline{https://wp.ufpel.edu.br/ermacrs23}}}.
      
    \bibitem[STROGATZ,  2015]{strogatz2015}Strogatz, S. H. {\bf Nonlinear Dynamics and Chaos: With Applications to Physics, Biology, Chemistry, and Engineering},  2nd ed. Philadelphia: Westview Press, 2015. ISBN: \texttt{978-0813349107}.
      
    \bibitem[VIANA, 2019]{Viana} Viana, D. S.  ``Condições de otimalidade para otimização cônica''. Tese de doutorado. Instituto de Matemática e estatística, Universidade de São Paulo, Programa de Matemática Aplicada, 2019.
  \end{thebibliography}
\end{document}
